\documentclass[10pt,a4paper]{amsart}
\usepackage[margin=19mm]{geometry}
\usepackage{amsmath,amssymb,mathtools}
\usepackage[scr=rsfs,cal=euler]{mathalfa}
\usepackage{xfrac}
\usepackage[unicode,hidelinks,bookmarks=false]{hyperref}
\newcommand{\ie}{i.e.\ }
\newcommand{\cf}{cf.\ }
\newcommand{\resp}{respectively\ }
\newcommand{\Subsection}{\subsection}
\providecommand{\nobreakdash}{\nobreak\hbox{-}}
\setlength{\emergencystretch}{2em}
\multlinegap0pt
\usepackage{bm}
\usepackage{bbm}
\usepackage{dsfont}
\usepackage{xspace}
\usepackage{cancel}
\usepackage{mathtools}
\usepackage{amsthm}
\makeatletter
\@ifundefined{theorem}{\newtheorem{theorem}{Theorem}}{}
\@ifundefined{lemma}{\newtheorem{lemma}[theorem]{Lemma}}{}
\makeatother
\makeatletter
\def\R{{\mathbb R}}
\def \SV{{SV}_{\delta}}
\expandafter\ifx\csname customproof\endcsname\relax
\newenvironment{customproof}[1][]
{\par\noindent{\textbf{Proof of #1:}}\quad}
\fi
\def \h{{\rm ht}}
\expandafter\ifx\csname pf\endcsname\relax
\newenvironment{pf}{\noindent {\bf Proof:}}{$\Box$ \vspace{2 ex}}
\fi
\expandafter\ifx\csname proof\endcsname\relax
\newenvironment{proof}{\noindent {\bf Proof:}}{$\Box$ \vspace{2 ex}}
\fi
\makeatother
\title[Boundedness of average rank of elliptic curves ordered by th]{Boundedness of average rank of elliptic curves ordered by the coefficients}
\author{Fatemehzahra Janbazi}
\date{Source: arXiv:2506.07089v1 (CC BY 4.0)}
\begin{document}
\maketitle

\noindent\emph{Source and licence.} Boundedness of average rank of elliptic curves ordered by the coefficients. Version arXiv:2506.07089v1 (CC BY 4.0), distributed under the Creative Commons Attribution 4.0 International licence, as recorded in the arXiv licence metadata for this exact version and shown on the abstract page https://arxiv.org/abs/2506.07089v1. Full rights notice: licenses/arxiv-2506.07089-CC-BY-4.0.txt.\par\smallskip

\noindent\textit{Editorial scope.} Counted unit: the Lemma whose source label is \texttt{help fglueing lemma} in the article (source0.tex lines 1652--1668, source label \texttt{help fglueing lemma}), delivered together with the article's own complete original proof of it (source0.tex lines 1671--1761, one \texttt{proof} environment of 91 lines). No mathematics is written, repaired or re-derived: statement and proof are the article's own text line for line. Required context retained uncounted: R0Hp (lines 837--839); RoHn (lines 850--852). Every definition, display and label of those lines is retained unchanged. Omitted from this package and not redistributed: 1--836, 840--849, 853--1651, 1669--1670, 1762--2288. Nothing inside a retained excerpt is omitted or altered: every retained body is delivered exactly as the article's lines are written, so no omission receipt of any kind accompanies this package. Citations: the retained excerpts carry no citation command at all, so no bibliography entry of the article is transcribed and no reference is redistributed. Reference adaptation: none was needed and none was made; every reference inside the delivered unit resolves inside it, so no unresolved reference or citation is produced. Printed slips kept literal: no printed word, symbol, hypothesis or number of the counted statement or of its proof was altered. Layout and class: the article's own class is \texttt{article} with packages including fontenc, amsfonts, amsmath, amssymb, latexsym, mathrsfs, comment, url, xy, mathtools, float, enumitem, xcolor, environ; the wrapper is the lane's pinned \texttt{amsart} editorial wrapper at 10pt on A4, which restates the theorem-like environments the retained text needs and adds the article's own macro definitions that the retained text needs (\texttt{\textbackslash R}, \texttt{\textbackslash SV}, \texttt{\textbackslash h}), with extra wrapper packages (bm, bbm, dsfont, xspace, cancel, mathtools). The delivered numbering is the unit's own. Assets: no retained span contains a figure environment, a graphics-inclusion command, a PDF-page inclusion, a table or a TikZ picture; no image file is copied into this package, the package declares no asset dependency and no image file is redistributed with it. Licence: version arXiv:2506.07089v1 of this article is distributed under the Creative Commons Attribution 4.0 International licence, as recorded in the arXiv licence metadata for this exact version and shown on the abstract page https://arxiv.org/abs/2506.07089v1. A full rights notice is delivered at licenses/arxiv-2506.07089-CC-BY-4.0.txt. Characters: every retained excerpt is pure ASCII, so the delivered engine, XeLaTeX with its default Latin Modern text font, reports no missing character.
    \begin{equation}\label{R0Hp}
        X - \frac{H^2}{48a^2} > \frac{4C}{3} \left| \frac{aX}{H} \right|,
    \end{equation}

    \begin{equation}\label{RoHn}
        R_0(a, b, c) = \sqrt{\left| \frac{48a^2 H}{27} \left( X + \frac{H^2}{48a^2} \right) \right|}.
    \end{equation}

\begin{lemma}\label{help fglueing lemma} 
Define $\SV'$ in the semi-invariant space by:
\begin{equation*}
\SV' \vcentcolon= \Big\{(a,b,c,R,I) : (\lfloor a \rfloor, \lfloor b \rfloor, \lfloor c \rfloor) \in B'_{X,\delta} \Big\}. 
\end{equation*} 
For any $(\lambda, t)$ and a given semi-form $f_s = (a, b, c, R, I) \in V_{\R}^s$, define the action $(\lambda, t) \cdot f_s$ by 
\begin{equation*} 
(\lambda, t) \cdot f_s \vcentcolon= \Big(\frac{at^4}{\lambda}, \frac{bt^2}{\lambda}, \frac{c}{\lambda}, \frac{Rt^6}{\lambda^3}, \frac{I}{\lambda^2}\Big).
\end{equation*}
Then we have
\begin{align*} 
&\sum_{(a,b,c) \in B'_{X,\delta}} \frac{1}{|8a^2 \cdot 12a|} \int_R \int_I \psi_{n} \left( \frac{at^4}{\lambda}, \frac{bt^2}{\lambda}, \frac{c}{\lambda}, \frac{Rt^6}{\lambda^3}, \frac{I}{\lambda^2} \right) \h_X'(a,b,c,R,I) \, dR \, dI 
= \\
&\quad\int_{f_s \in \SV'} \frac{1}{|8a^2 \cdot 12a|} , \psi_{n}\left((\lambda, t) \cdot f_s\right) \h_X'(f_s) \, df_s +
O(X^{3/2 +200\delta}) \end{align*} 
where $\h_X'(f_s)$ is defined as before.
\end{lemma}

\begin{proof}  Define the function $g(a, b, c)$ as  
\begin{equation*}
    g(a, b, c) \vcentcolon= \frac{R_0(a, b, c) X}{a^2 H},
\end{equation*}  
where $R_0(a,b,c)$ is given in equations \eqref{R0Hp} and \eqref{RoHn}.  
Next, define the function $F(a, b, c, R, I)$ as
\begin{equation*}
    F(a, b, c, R, I) \vcentcolon= \psi_{n} \Big( \frac{at^4}{\lambda}, \frac{bt^2}{\lambda}, \frac{c}{\lambda}, \frac{RR_0 t^6}{\lambda^3}, \frac{\frac{4aCXI}{3H} + \frac{4aC\Gamma}{3H}}{\lambda^2} \Big),
\end{equation*}
where $\Gamma$ is $\Gamma(a,b,c,R)$, as defined previously.

Then, after a suitable change of variables, the difference of the integrals in the lemma simplifies to
\begin{align*} 
\int_{I} \int_{R}  \chi_{(-1,1)}(R) \chi_{(-1,1)}(I) \Big( \sum_{(a,b,c)\in B'_{X,\delta}} \int_{a}^{a+1} \int_{b}^{b+1} \int_{c}^{c+1} & \Big( g(a,b,c) F(a, b, c, R, I) \\
&\quad - g(a',b',c') F(a', b', c', R, I) \Big) \, d\epsilon_a \, d\epsilon_b \, d\epsilon_c \Big)\, dR \, dI,
\end{align*}
where we define $a' = a + \epsilon_a$, $b' = b + \epsilon_b$, and $c' = c + \epsilon_c$ for small $\epsilon_a$, $\epsilon_b$, and $\epsilon_c$.
\noindent To complete the proof, we need to show that the difference 
\[
\left| F(a, b, c, R, I) g(a, b, c) - F(a', b', c', R, I) g(a', b', c') \right|
\]
is small. Thus, we need to bound both $|g(a, b, c) - g(a', b', c')|$ and $|g(a,b,c)\big(F(a, b, c, R, I) - F(a', b', c', R, I)\big)|$. For clarity, we will explicitly demonstrate the bound for $F(a, b, c, R, I)$, as the argument for $g(a, b, c)$ follows analogously.

Since $g(a, b, c)$ is bounded by $O(X^{1/2 + 2 \delta_a + \delta_H/2})$ for all $(a, b, c) \in B'_{X, \delta}$, it suffices to focus on bounding the difference between the values of $F(a, b, c, R, I)$ and $F(a', b', c', R, I)$, where:
\begin{equation*}
F(a, b, c, R, I) \vcentcolon= \psi_{n} \Big( \frac{a t^4}{\lambda}, \frac{b t^2}{\lambda}, \frac{c}{\lambda}, \frac{R R_0 t^6}{\lambda^3}, \frac{\frac{4aCXI}{3H} +  \frac{27 R_0^2R^2}{48 a^2 H} + \frac{H^2}{48 a^2}}{\lambda^2}\Big),
\end{equation*}
and
\begin{equation*}
F(a', b', c', R, I) \vcentcolon= \psi_{n} \Big( \frac{a' t^4}{\lambda}, \frac{b' t^2}{\lambda}, \frac{c'}{\lambda}, \frac{R R_0' t^6}{\lambda^3}, \frac{R R_0 t^6}{\lambda^3}, \frac{\frac{4a'CXI}{3H'} +  \frac{27 R_0'^2R^2}{48 (a')^2 H'} + \frac{(H')^2}{48 (a')^2}}{\lambda^2} \Big).
\end{equation*}
By the Mean Value Theorem, it suffices to show that each component of $\psi_{n}$ remains close when evaluated at $(a, b, c)$ and $(a', b', c')$, as $\psi_{n}$ is both bounded and smooth. This requires demonstrating that differences such as

\begin{equation*}
 \lambda^{-1}|a t^4 - a' t^4|, \quad \lambda^{-1}| b t^2 -b' t^2|, \quad \lambda^{-1}|c-c'|, \quad  \lambda^{-3}t^6|R_0-R_0' |,
\end{equation*}
and 
\begin{equation*}
\Big| \frac{ a }{H } - \frac{a' }{H' } \Big|, \quad 
\lambda^{-2}\left|\Big(\frac{27 R_0^2R^2}{48 a^2 H} + \frac{H^2}{48 a^2}\Big)-\Big(\frac{27 R_0'^2R^2}{48 (a')^2 H'} + \frac{(H')^2}{48 (a')^2}\Big) \right|,
\end{equation*}
are all bounded by $O(X^{-1/2 + \delta'})$ for a sufficiently small $\delta' > 0$. We analyze each term separately:

\begin{itemize}
    \item Bounding $\lambda^{-1}|a t^4 - a' t^4|$:
    \begin{equation*}
    \Big| \frac{a t^4}{\lambda} - \frac{a' t^4}{\lambda} \Big| = \frac{|\epsilon_a| t^4}{\lambda} \ll X^{-1/2 + 4\delta_t}.
    \end{equation*}
Similar argument applies to $\lambda^{-1}| b t^2 -b' t^2|$ and $\lambda^{-1}|c-c'|$, yielding $O(X^{-1/2 + 4\delta_t})$ for each.

\item Bounding $|R_0-R_0'|t^6\lambda^{-3}$ :
we aim to show that $R_0$ and $R_0'$ are sufficiently close for any $(a, b, c) \in B'_{X, \delta}$. Given the polynomial representation \( H(x,y,z) = 3y^2 - 8xz \), the Mean Value Theorem implies
\begin{equation*}
|H(a', b', c') - H(a, b, c)| = O(X^{1/2}).
\end{equation*}
Since $|H(a, b, c)|$ is greater than $X^{1 - \delta_H}$, this ensures that $H$ and $H'$ maintain the same sign.

To make the analysis concrete, we first consider the case where \( H(a, b, c) < 0 \); the case for \( H(a, b, c) > 0 \) will follow similarly. Based on the definition of \( R_0 \), we need to show that the difference
\begin{equation*}
\Big| |a| |H|^{1/2} \sqrt{X + \frac{H^2}{48a^2}} - |a'| |H'|^{1/2} \sqrt{X + \frac{(H')^2}{48(a')^2}} \Big|
\end{equation*}
is sufficiently small. Since we have upper and lower bounds on $a$ and $H$, we can decompose this expression into three independent terms and bound each difference separately.

Firstly, we consider the term
\begin{equation*}
\Big| |a| |H|^{1/2} \sqrt{X + \frac{H^2}{48a^2}} - |a'| |H|^{1/2} \sqrt{X + \frac{H^2}{48a^2}} \Big|.
\end{equation*}
Here, the difference is $O(X^{1/2} \cdot X^{1/2 + \delta_a})$ because both $H$ and $a$ are bounded above and below, with $|a - a'| = O(1)$ by assumption. This term contributes an error of order $O(X^{1 + \delta_a})$.

Secondly, we examine
\begin{equation*}
\Big| |a'| |H|^{1/2} \sqrt{X + \frac{H^2}{48a^2}} - |a'| |H'|^{1/2} \sqrt{X + \frac{H^2}{48a^2}} \Big|.
\end{equation*}
This term is $O(X^{1 + \delta_H/2 + 2\delta_a})$ because $H$ and $H'$ have the same sign, and we have the constraints $|H - H'| = O(X^{1/2})$ and $|H| > X^{1 - \delta_H}$. This bound follows directly from these constraints.

Finally, we consider the term
\begin{equation*}
\Big| |a'| |H'|^{1/2} \sqrt{X + \frac{H^2}{48a^2}} - |a'| |H'|^{1/2} \sqrt{X + \frac{(H')^2}{48(a')^2}} \Big|.
\end{equation*}
Our goal is to show that this difference is also small. Using the Mean Value Theorem, we know that $|(H/a)^2 - (H'/a')^2| = O(X^{1/2 + 3\delta_a})$, which implies that this term is bounded by $O(X^{1 + 3\delta_a})$.

In summary, each of the three terms is bounded, resulting in a total error of $O(X^{-1/2 + 3\delta_a+\delta_H+6\delta_t})$ for $|(R_0-R_0')t^6\lambda^{-3}|$.
\item Bounding $ \left| \frac{a}{H} - \frac{a'}{H'} \right|$: Since the difference between $H$ and $H'$ is $O(X^{1/2})$ and $H$ is bounded from below, we obtain the bound $O(X^{-1+2\delta_H})$ for this term.
    
\item Bounding $ \left| \frac{27 R_0^2}{48 a^2 H} - \frac{27 (R_0')^2}{48 (a')^2 H'} \right|$: We use the definition along with the fact that $a'$ and $H'$ have the same sign as $a$ and $H$. This reduces to bounding $ \Big| \Big(\frac{H}{a}\Big)^2 - \Big(\frac{H'}{a'}\Big)^2 \Big|$, which we have shown to be as small as $O(X^{1/2+3\delta_a})$. Consequently:
    \begin{equation*}
        \lambda^{-2} \Big| \Big( \frac{27 R_0^2 R^2}{48 a^2 H} + \frac{H^2}{48 a^2} \Big) - \Big( \frac{27 (R_0')^2 R^2}{48 (a')^2 H'} + \frac{(H')^2}{48 (a')^2} \Big) \Big| = O(X^{-1/2+3\delta_a}).
    \end{equation*}
\end{itemize}
Adding all the errors will complete the proof.
\end{proof}

\end{document}
